\begin{center}
    {\LARGE \textbf{2005无机化学}} \\[2ex]
    {\large 时量180分钟 \quad 满分150分}
\end{center}

\vspace{0.5cm}
\noindent\textbf{注意事项:}
\begin{enumerate}[label=\arabic*., parsep=0pt, itemsep=0pt]
    \item 答题前,考生先将自己的姓名、学号、考场号、座位号填写在答题卡上,将条形码准确粘贴在考生信息条形码粘贴区。
    \item 考试结束后,将本试卷与答题纸一并收回。
    \item 回答各个问题时,将答案写在答题纸上,写在本试卷上无效。
\end{enumerate}

\vspace{0.5cm}
\noindent\textbf{一、按要求回答下列问题。(15 pts.)}

\begin{enumerate}[label=\arabic*., start=1, itemsep=1.5ex]
    \item 将\ce{KClO}溶液滴入\ce{KI}溶液中。
    \item 向\ce{H2O2}溶液中滴加酸性\ce{KMnO4}溶液。
    \item 将白磷固体投入溴水中。
    \item 用盐酸处理\ce{Co2O3}。
    \item 向碘水溶液中滴加\ce{NaHSO3}溶液。
\end{enumerate}

\vspace{0.5cm}
\noindent\textbf{二、给出实验现象和对应的化学反应方程式。(15 pts.)}

\begin{enumerate}[label=\arabic*., start=1, itemsep=1.5ex]
    \item 向酸性\ce{KIO3}和淀粉混合溶液中滴加\ce{Na2SO3}溶液。
    \item 将\ce{SnCl2}溶液滴入\ce{HgCl2}溶液。
    \item 向\ce{K2Cr2O7}溶液中通入过量\ce{SO2}。
    \item 将氨水滴加到\ce{NiSO4}溶液中。
    \item 把\ce{MnSO4}、\ce{H2O2}及\ce{NaOH}三种溶液混合。
\end{enumerate}

\vspace{0.5cm}
\noindent\textbf{三、解释下列实验现象。(20 pts.)}

\begin{enumerate}[label=\arabic*., start=1, itemsep=1.5ex]
    \item 把少量碘水滴入\ce{FeSO4}溶液，碘水不褪色；再加入适量\ce{NaHCO3}溶液后，则碘水褪色。
    \item 混合\ce{FeCl3}和\ce{KSCN}溶液后，溶液为红色；再加入适量\ce{NH4F}或\ce{Zn}粉，溶液褪色。
    \item 铜片在氨水中腐蚀的速度比在\ce{NaCl}溶液中快。
    \item \ce{Na3PO4}溶液中滴加\ce{AgNO3}溶液时生成黄色沉淀，但向\ce{NaPO3}溶液中滴加\ce{AgNO3}溶液时却生成白色沉淀。
    \item 在煤气灯上加热\ce{KNO3}晶体时没有棕色气体生成；若在\ce{KNO3}晶体混有\ce{CuSO4}时则有棕色气体生成。
\end{enumerate}

\vspace{0.5cm}
\noindent\textbf{四、简要回答下列各题。}

\begin{enumerate}[label=\arabic*., start=1, itemsep=1.5ex]
    \item 给出下列分子或离子的几何构型和中心原子的杂化类型。(7 pts.)
    
    {\centering
    \ce{SOCl2}，\ce{SO3}，\ce{NH3}，\ce{O3}，\ce{ICl4^{-}}，\ce{AuCl4^{-}}，\ce{XeF4} \par
    }

    \item 指出下列分子是否含有特殊键型（离域大$\Pi$键，d-p$\pi$配键）并给出是哪种特殊键型。(7 pts.)
    
    {\centering
    \ce{O3}，\ce{SO3}，\ce{SO2Cl2}，\ce{HClO3}，\ce{HNO2}，\ce{HNO3}，\ce{H3PO4} \par
    }

    \item 下列配合物是否有旋光异构体，如果有请画出来。(4 pts.)
    
    {\centering
    \ce{[PtCl4(NH3)2]}，\ce{[CoCl2(en)(NH3)2]^{+}}，\ce{[CoCl2(en)2]^{-}}，\ce{[PtCl4(en)]} \par
    }

    \item 已知：\\
    \ce{Fe^{3+} + e^{-} = Fe^{2+}} \quad $E^\ominus = 0.77\ \text{V}$ \\
    \ce{Fe(CN)6^{3-} + e^{-} = Fe(CN)6^{4-}} \quad $E^\ominus = 0.36\ \text{V}$

    试比较\ce{Fe(CN)6^{3-}}与\ce{Fe(CN)6^{4-}}哪个稳定常数更大？并予以说明。(4 pts.)

    \item 比较配合物的稳定性并说明原因。(4 pts.)
    \begin{enumerate}[label=(\arabic*), itemsep=1ex]
        \item \ce{Cu(NH3)4^{2+}}与\ce{Zn(NH3)4^{2+}}
        \item \ce{ZnI4^{2-}}与\ce{HgI2^{2-}}
    \end{enumerate}

    \item 给出元素的价电子结构：\ce{Cr}，\ce{Cu}，\ce{Ru}，\ce{Pd} (4 pts.)
    
    \item 已知\ce{NO2}的键角（$132^\circ$）比预想的大得多，请给出合理的解释。(4 pts.)
    
    \item 请解释原因：(6 pts.)
    \begin{enumerate}[label=(\arabic*), itemsep=1ex]
        \item 分子的极性 \quad \ce{NH3} $>$ \ce{NF3}
        \item 热稳定性 \quad \ce{K2CO3} $>$ \ce{MgCO3}
        \item 熔沸点 \quad \ce{NH3} $<$ \ce{HF}
    \end{enumerate}

    \item 已知物质的颜色：\ce{VF3} 黄绿色，\ce{VCl3} 紫红色，\ce{VBr3} 棕灰色。请解释颜色发生变化的原因。(6 pts.)
    
    \item 已知下列物质的熔点，请对熔点的变化加以解释。(6 pts.)
    
    \vspace{1ex}
    \begin{tabular}{c c c c c}
    & \ce{TiF4} & \ce{TiCl4} & \ce{TiBr4} & \ce{TiI4} \\
    熔点 & $284^\circ\text{C}$ & $-24^\circ\text{C}$ & $38^\circ\text{C}$ & $150^\circ\text{C}$
    \end{tabular}
    \vspace{1ex}

\end{enumerate}

\vspace{0.5cm}
\noindent\textbf{五、制备与分离。}

\begin{enumerate}[label=\arabic*., start=1, itemsep=1.5ex]
    \item 用反应方程式（不用配平）和简要的文字表示下列转化过程：\ce{Fe -> FeSO4 * 7H2O -> FeCl3 * 6H2O -> K3[Fe(C2O4)3] * 3H2O} (6 pts.)
    \item 用反应方程式表示：由\ce{Cu}为主要原料制备\ce{CuI}。(5 pts.)
    \item \ce{CuCl}、\ce{AgCl}、\ce{Hg2Cl2}都是难溶于水的白色粉末，试区别这三种物质。(6 pts.)
    \item 用二种方法鉴别下列各对物质：(6 pts.)
    \begin{enumerate}[label=(\arabic*), itemsep=1ex]
        \item \ce{NH4NO3}和\ce{NH4I}；
        \item \ce{NaNO3}和\ce{NaNO2}；
        \item \ce{NaCl}和\ce{NaI}
    \end{enumerate}
    \item 设计方案分离\ce{Al^{3+}}、\ce{Cr^{3+}}、\ce{Fe^{3+}}、\ce{Co^{2+}}、\ce{Zn^{2+}}。(5 pts.)
\end{enumerate}

\vspace{0.5cm}
\noindent\textbf{六、推断题。(20 pts.)}

\begin{enumerate}[label=\arabic*., start=1, itemsep=1.5ex]
    \item 将无色钠盐溶于水得无色溶液\ce{A}，用\ce{pH}试纸检验知\ce{A}显酸性。向\ce{A}中滴加\ce{KMnO4}溶液，则紫红色褪去，说明\ce{A}被氧化为\ce{B}，向\ce{B}中加入\ce{BaCl2}溶液得不溶于强酸的白色沉淀\ce{C}。向\ce{A}中加入稀盐酸有无色气体\ce{D}放出，将\ce{D}通入\ce{KMnO4}溶液则又得到无色的\ce{B}。向含有淀粉的\ce{KIO3}溶液中滴加少许\ce{A}则溶液立即变蓝，说明有\ce{E}生成，\ce{A}过量时蓝色消失得无色溶液，试给出\ce{A,B,C,D,E}的化学式。
    \item 绿色固体\ce{A}溶于水后通入\ce{CO2}即得棕黑色固体\ce{B}和紫红色溶液\ce{C}。用酸性\ce{H2O2}溶液处理\ce{B}则\ce{B}溶解得到近无色溶液\ce{D}。向\ce{D}中加入\ce{NaOH}溶液得到白色沉淀\ce{E}，\ce{E}暴露在空气中缓慢转化为\ce{B}。\ce{B}与\ce{KOH}、\ce{KClO3}共熔得到的产物含有\ce{A}。溶液\ce{C}与酸性\ce{H2O2}溶液混合也生成\ce{D}。试判断\ce{A,B,C,D,E}各为何物质？
    
\end{enumerate}